% ============================================================
% Round 02: Diffusion / Flow Minimal Foundations for Post-Training
% ============================================================

\subsection{Round 02：Diffusion / Flow 后训练最小基础}

\subsubsection{本轮问题}

\begin{conceptbox}
\textbf{学习目标：} 只补后训练需要的 Diffusion / Flow 基础：trajectory 是什么、action 是什么、logprob 从哪里来、最终图像 reward 如何作用到中间生成步骤。\\
\textbf{主线位置：}
\[
\text{图像编辑任务本体}
\to
\text{Diffusion / Flow 后训练基础}
\to
\text{Reward / Preference}
\to
\text{Diffusion-DPO / DDPO / Flow-GRPO}
\]
\end{conceptbox}

\subsubsection{为什么要先补 Diffusion / Flow}

文本侧后训练里，policy 很清楚：

\[
\pi_\theta(y_t \mid x,y_{<t})
\]

每一步 action 是下一个 token，logprob 也很自然：

\[
\log \pi_\theta(y_t \mid x,y_{<t})
\]

但图像生成 / 编辑模型不是逐 token 生成，而是从噪声 latent 逐步变成图像：

\[
z_T
\to
z_{T-1}
\to
\cdots
\to
z_0
\to
I
\]

所以图像后训练首先要回答：

\[
\text{这条生成轨迹里，哪一步算 action？}
\]

\[
\text{每一步有没有概率？能不能算 logprob？}
\]

\[
\text{最终图像的 reward 怎么分配给中间步骤？}
\]

\begin{summarybox}
\textbf{一句话：} 文本后训练优化 token distribution；图像后训练优化 denoising / flow trajectory。
\end{summarybox}

\subsubsection{Diffusion 的最小图景}

Diffusion model 有两个过程：

\[
\text{forward noising}
\qquad
\text{reverse denoising}
\]

\paragraph{Forward noising}

训练时，把真实图像 latent $z_0$ 逐步加噪：

\[
z_0
\to
z_1
\to
\cdots
\to
z_T
\]

其中 $z_T$ 接近纯高斯噪声。

常见写法是：

\[
q(z_t \mid z_0)
=
\mathcal{N}
\left(
\sqrt{\bar{\alpha}_t}z_0,\;
(1-\bar{\alpha}_t)I
\right)
\]

直觉是：$t$ 越大，图像信息越少，噪声越多。

\paragraph{Reverse denoising}

推理时，从噪声 $z_T$ 出发，一步步去噪：

\[
z_T
\to
z_{T-1}
\to
\cdots
\to
z_0
\]

模型学习的是反向转移：

\[
p_\theta(z_{t-1}\mid z_t,c)
\]

其中 $c$ 是条件。对图像编辑来说，条件通常包括：

\[
c = (I_\text{src}, \text{edit instruction})
\]

\begin{intubox}
\textbf{直觉：} Diffusion 的每一步都在问：当前这个 noisy latent 应该怎么变得更干净一点，同时符合原图和编辑指令？
\end{intubox}

\subsubsection{Diffusion 里的 Policy、Action、Trajectory}

从 RL 角度看，reverse denoising 可以看成一个多步决策过程。

\begin{center}
{\small
\begin{tabular}{p{3cm}p{9.8cm}}
\hline
\textbf{RL 概念} & \textbf{Diffusion / Image Editing 中的含义} \\
\hline
State $s_t$
&
当前 latent $z_t$，时间步 $t$，原图 $I_\text{src}$，编辑指令 $c$
\\
Action $a_t$
&
采样下一个 latent $z_{t-1}$，或模型预测的噪声 / velocity / score
\\
Policy $\pi_\theta$
&
反向去噪模型 $p_\theta(z_{t-1}\mid z_t,I_\text{src},c,t)$
\\
Trajectory
&
$z_T,z_{T-1},\ldots,z_0$ 这整条去噪路径
\\
Reward
&
最终图像 $I_\text{edit}$ 的编辑质量、保真度、美学、人类偏好分数
\\
\hline
\end{tabular}
}
\end{center}

如果反向转移是随机高斯：

\[
z_{t-1}
\sim
\mathcal{N}
\left(
\mu_\theta(z_t,t,c),\;
\sigma_t^2 I
\right)
\]

那么每一步都有 logprob：

\[
\log p_\theta(z_{t-1}\mid z_t,c,t)
\]

这就能接上 PPO / DDPO / GRPO：

\[
\rho_t(\theta)
=
\frac{
p_\theta(z_{t-1}\mid z_t,c,t)
}{
p_{\theta_\text{old}}(z_{t-1}\mid z_t,c,t)
}
\]

\begin{summarybox}
\textbf{Diffusion RL 的关键：} 只要每一步反向采样有概率密度，就能像文本 PPO 那样计算 old logprob、new logprob 和 ratio。
\end{summarybox}

\subsubsection{最终 Reward 怎么给中间步骤}

图像 reward 通常只在最后才知道：

\[
I_\text{edit}
=
\operatorname{Decode}(z_0)
\]

\[
R
=
R(I_\text{src},c,I_\text{edit})
\]

这个 reward 是整张最终图像的分数，而不是每一步 $z_t$ 都有分数。

所以最简单的做法是：整条 trajectory 共享同一个 reward 或 advantage：

\[
A_t \approx A
\]

于是所有去噪步骤都朝同一个方向更新：

\[
\sum_t
\log p_\theta(z_{t-1}\mid z_t,c,t)
\cdot A
\]

\begin{warnbox}
\textbf{信用分配难点：} 如果最终图像不好，到底是哪一步去噪造成的？早期构图错了？中期语义错了？后期细节坏了？只用最终 reward 很难精确回答。
\end{warnbox}

这就是图像后训练里很多方法要研究 timestep weighting、early-step optimization、trajectory reconstruction、process reward 的原因。

\subsubsection{Flow Matching 的最小图景}

Flow Matching / Rectified Flow 的思路和 diffusion 不一样。

它通常构造一条从真实图像 $x_0$ 到噪声 $x_1$ 的连续路径：

\[
x_t = (1-t)x_0 + t x_1,
\qquad
t\in[0,1]
\]

其中：

\[
t=0 \Rightarrow x_t=x_0
\]

\[
t=1 \Rightarrow x_t=x_1
\]

模型学习速度场：

\[
v_\theta(x_t,t,c)
\]

训练目标是让预测速度接近真实路径速度：

\[
\mathcal{L}_\text{FM}
=
\mathbb{E}
\left[
\left\|
v_\theta(x_t,t,c)
-
(x_1-x_0)
\right\|^2
\right]
\]

推理时，从噪声出发，沿 ODE 积分回图像：

\[
\frac{dx_t}{dt}
=
v_\theta(x_t,t,c)
\]

离散写成：

\[
x_{t-\Delta t}
=
x_t
-
v_\theta(x_t,t,c)\Delta t
\]

\begin{intubox}
\textbf{直觉：} Diffusion 更像“逐步去噪”；Flow Matching 更像“学习一个速度场，让噪声沿着流动路径变成图像”。
\end{intubox}

\subsubsection{Flow 的 RL 难点：ODE 没有 Logprob}

Flow Matching 默认推理是 ODE：

\[
x_{t-\Delta t}
=
f_\theta(x_t,t,c)
\]

给定当前状态 $x_t$、时间 $t$ 和条件 $c$，下一步是确定的。

这对生成很方便，但对 RL 很麻烦：

\begin{itemize}[leftmargin=1.5em]
  \item 没有随机性，探索弱；
  \item 没有正常的概率密度，难以计算 $\log p_\theta(x_{t-\Delta t}\mid x_t,c,t)$；
  \item 没有 logprob，就很难计算 PPO / GRPO 的 ratio。
\end{itemize}

\begin{warnbox}
\textbf{核心矛盾：} 文本 PPO / diffusion RL 需要 logprob；deterministic flow ODE 默认不给你 logprob。
\end{warnbox}

所以 Flow-GRPO 这类方法会把 ODE 改造成带随机性的 SDE：

\[
x_{t-\Delta t}
=
\mu_\theta(x_t,t,c)
+
\sigma_t\sqrt{\Delta t}\epsilon,
\qquad
\epsilon\sim\mathcal{N}(0,I)
\]

这样每一步转移就变成高斯：

\[
p_\theta(x_{t-\Delta t}\mid x_t,c,t)
=
\mathcal{N}
\left(
x_{t-\Delta t};
\mu_\theta(x_t,t,c),
\sigma_t^2\Delta t I
\right)
\]

于是可以计算：

\[
\log p_\theta(x_{t-\Delta t}\mid x_t,c,t)
\]

也就可以做 GRPO ratio。

\begin{summarybox}
\textbf{Flow-GRPO 的最小动机：} 把 deterministic ODE 变成 stochastic SDE，让 flow trajectory 既能探索，也能计算 logprob。
\end{summarybox}

\subsubsection{Diffusion 和 Flow 的后训练差别}

\begin{center}
{\small
\begin{tabular}{p{3cm}p{4.8cm}p{4.8cm}}
\hline
 & \textbf{Diffusion} & \textbf{Flow Matching} \\
\hline
生成过程
&
随机去噪 Markov chain
&
通常是确定性 ODE flow
\\
每步概率
&
通常天然有 $p_\theta(z_{t-1}\mid z_t)$
&
ODE 默认没有，需要 SDE 化或其他技巧
\\
RL 接入
&
较自然，可算 denoising step logprob
&
要先引入随机性和可计算 logprob
\\
trajectory
&
denoising trajectory
&
flow / ODE / SDE trajectory
\\
典型方法
&
DDPO, Diffusion-DPO
&
Flow-GRPO, LeapAlign
\\
\hline
\end{tabular}
}
\end{center}

\subsubsection{图像编辑条件如何进入模型}

图像编辑不是只给文本条件 $c$。更准确是：

\[
c_\text{edit}
=
(I_\text{src}, \text{instruction})
\]

不同模型会用不同方式注入原图：

\begin{itemize}[leftmargin=1.5em]
  \item 把原图编码成 latent，与 noisy latent 拼接；
  \item 用 VAE encoder 提供外观控制；
  \item 用 VLM / image encoder 提供语义控制；
  \item 用 mask 指定可编辑区域；
  \item 用 cross-attention 注入文本编辑指令。
\end{itemize}

因此图像编辑的反向转移更像：

\[
p_\theta(z_{t-1}
\mid
z_t,\;
I_\text{src},\;
c,\;
t)
\]

\begin{intubox}
\textbf{直觉：} text-to-image 只问“按文本生成什么”；image editing 还要问“基于这张原图，哪些内容该继承，哪些内容该改变”。
\end{intubox}

\subsubsection{本轮最小结论}

\begin{summarybox}
\textbf{Diffusion / Flow 后训练最小基础：}\\
Diffusion 生成是一条随机去噪轨迹：
\[
z_T\to z_{T-1}\to\cdots\to z_0
\]
每一步通常有 logprob，因此能自然接 PPO / DDPO / GRPO。\\
Flow Matching 生成通常是确定性 ODE：
\[
\frac{dx_t}{dt}=v_\theta(x_t,t,c)
\]
默认没有 logprob，所以做 RL 时常需要 SDE 化，引入随机性和可计算转移概率。\\
图像编辑后训练的核心难点是：最终图像 reward 如何稳定地作用到 denoising / flow trajectory 的中间步骤。
\end{summarybox}

\subsubsection{本轮检查题}

\begin{enumerate}[leftmargin=1.5em]
  \item Diffusion 的 forward noising 和 reverse denoising 分别是什么？
  \item 在 diffusion RL 里，state、action、trajectory、reward 分别对应什么？
  \item 为什么 diffusion 的每步随机转移更容易接 PPO / GRPO？
  \item Flow Matching 的 ODE 为什么不能直接算 logprob？
  \item Flow-GRPO 为什么要把 ODE 改成 SDE？
\end{enumerate}
